\documentclass[times]{san}

\bibliografia{bibliografia.bib}

\uczelnia{Społeczna Akademia Nauk}
\przedmiot{Bazy danych i aplikacje}
\autor{Adrian}
\nralbumu{119205}
\grupa{4}
\tytul{System zgłaszania artykułów naukowych}
\prowadzacy{mgr inż. Krystian Gumiński}
\miasto{Łódź}
\rok{2026}

\begin{document}

\maketitlepage
\spistresci

\section{Wstęp}

Niniejszy dokument zawiera szczegółową dokumentację projektu ``System zgłaszania artykułów naukowych''. Projekt został zrealizowany jako aplikacja webowa działająca w architekturze klient-serwer z wykorzystaniem nowoczesnych technologii webowych. Celem systemu jest umożliwienie autorów naukowych przesyłania artykułów, zaś redaktorom i recenzentom -- przydzielania i publikowania recenzji artykułów. Aplikacja wspiera zarządzanie procesem redakcyjnym wraz z systemem powiadomień.

\section{Cel projektu}

Głównym celem projektu jest stworzenie kompleksowego systemu informatycznego do zarządzania procesem zgłaszania i recenzowania artykułów naukowych. Specyficzne cele obejmują:

\begin{itemize}
    \item Umożliwienie autorom przesyłania artykułów w formacie cyfrowym
    \item Automatyczne przydzielanie artykułów recenzentom do oceny
    \item Wspieranie komunikacji między autorami a recenzentami poprzez system powiadomień
    \item Automatyzacja procesu redakcyjnego z podziałem ról (autor, recenzent, administrator)
    \item Zabezpieczenie danych użytkowników i materiałów poprzez systemy autentykacji i autoryzacji
    \item Skalowalna architektura gotowa do rozwijania w przyszłości
\end{itemize}

Wymagania funkcjonalne systemu obejmują:
\begin{itemize}
    \item Rejestracja i logowanie użytkowników
    \item Zarządzanie profilami użytkowników
    \item Przesyłanie artykułów z metadanymi
    \item Przydzielanie artykułów recenzentom
    \item Prowadzenie recenzji artykułów
    \item Podejmowanie decyzji redakcyjnych
    \item System powiadomień o zdarzeniach
    \item Panel administracyjny do zarządzania systemem
\end{itemize}

Wymagania niefunkcjonalne dotyczą wydajności, bezpieczeństwa, niezawodności i łatwości utrzymania kodu.

\section{Przegląd architektury}

\subsection{Ogólny model architektoniczny}

System został zbudowany na architekturze trójwarstwowej:

\begin{table}[h]
\centering
\caption{Warstwy architektury systemu}
\label{tab:architecture}
\begin{tabular}{|c|l|p{6cm}|}
\hline
Lp. & Warstwa & Opis \\
\hline
1   & Prezentacji & Frontend React aplikacja działająca w przeglądarce \\ 
2   & Logiki biznesowej & Backend ASP.NET Core z API REST \\
3   & Danych & Baza danych SQLite \\
\hline
\end{tabular}
\end{table}

\subsection{Podział komponentów}

System składa się z następujących głównych komponentów:

\subsubsection{Backend (ASP.NET Core)}

Backend stanowi serce aplikacji, odpowiadający za całą logikę biznesową:

\begin{itemize}
    \item \textbf{Serwer API REST} -- udostępnia endpointy dla frontendu
    \item \textbf{Autentykacja JWT} -- obsługuje logowanie użytkowników i wydawanie tokenów JWT
    \item \textbf{Entity Framework Core} -- mapowanie obiektowo-relacyjne dla bazy danych
    \item \textbf{Autoryzacja} -- system ról i uprawnień
    \item \textbf{Powiadomienia} -- system powiadomień dla użytkowników
    \item \textbf{Monitoring} -- monitoring systemu za pomocą Prometheus
\end{itemize}

\subsubsection{Frontend (React + Vite)}

Frontend zapewnia interaktywny interfejs użytkownika:

\begin{itemize}
    \item \textbf{Strona logowania} -- logowanie i rejestracja użytkowników
    \item \textbf{Panel autora} -- możliwość przesyłania nowych artykułów oraz przeglądania историї
    \item \textbf{Panel recenzenta} -- przeglądanie przydzielonych artykułów i prowadzenie recenzji
    \item \textbf{Panel administratora} -- zarządzanie systemem, użytkownikami i recenzjami
    \item \textbf{System powiadomień} -- wyświetlanie powiadomień o zdarzeniach
\end{itemize}

\subsubsection{Baza danych (SQLite)}

Baza danych przechowuje wszystkie dane systemu:

\begin{itemize}
    \item \textbf{Użytkownicy} -- dane profili użytkowników i dane logowania
    \item \textbf{Zgłoszenia} -- artykuły przesłane do systemu
    \item \textbf{Recenzje} -- oceny artykułów od recenzentów
    \item \textbf{Przydzielenia recenzji} -- mapowanie artykułów na recenzentów
    \item \textbf{Powiadomienia} -- historia powiadomień dla użytkowników
\end{itemize}

\section{Model danych}

\subsection{Diagramy encji}

System operuje na następujących głównych encjach:

\begin{table}[h]
\centering
\caption{Główne encje systemu}
\label{tab:entities}
\begin{tabular}{|c|l|p{5cm}|}
\hline
Lp. & Encja & Opis \\
\hline
1   & User & Użytkownik systemu \\
2   & Submission & Artykuł przesłany do systemu \\
3   & Review & Recenzja artykułu \\
4   & ReviewAssignment & Przydzielenie artykułu recenzentowi \\
5   & Notification & Powiadomienie dla użytkownika \\
6   & FileRecord & Rekord pliku przesłanego \\
\hline
\end{tabular}
\end{table}

\subsection{User (Użytkownik)}

Encja reprezentująca użytkownika systemu:

\begin{itemize}
    \item \textbf{Id} -- unikalny identyfikator (guid)
    \item \textbf{Email} -- adres email, unikalny
    \item \textbf{Name} -- imię i nazwisko
    \item \textbf{Role} -- rola w systemie: admin, author, reviewer
    \item \textbf{PasswordHash} -- zahashowane hasło (BCrypt)
\end{itemize}

\subsection{Submission (Zgłoszenie)}

Encja reprezentująca artykuł przesłany do systemu:

\begin{itemize}
    \item \textbf{Id} -- unikalny identyfikator
    \item \textbf{AuthorId} -- referencja do autora (User)
    \item \textbf{Title} -- tytuł artykułu
    \item \textbf{Abstract} -- streszczenie artykułu
    \item \textbf{Content} -- zawartość artykułu (tekst)
    \item \textbf{CreatedAt} -- data utworzenia
    \item \textbf{UpdatedAt} -- data ostatniej aktualizacji
    \item \textbf{Status} -- status: submitted, under-review, accepted, rejected
    \item \textbf{Decision} -- ostateczna decyzja redakcyjna
\end{itemize}

\subsection{Review (Recenzja)}

Encja reprezentująca recenzję artykułu:

\begin{itemize}
    \item \textbf{Id} -- unikalny identyfikator
    \item \textbf{SubmissionId} -- referencja do artykułu (Submission)
    \item \textbf{ReviewerId} -- referencja do recenzenta (User)
    \item \textbf{Content} -- treść recenzji
    \item \textbf{Recommendation} -- rekomendacja: accept, reject, major-revisions, minor-revisions
    \item \textbf{CreatedAt} -- data utworzenia recenzji
    \item \textbf{UpdatedAt} -- data ostatniej aktualizacji
\end{itemize}

\subsection{ReviewAssignment (Przydzielenie recenzji)}

Encja reprezentująca przydzielenie artykułu recenzentowi:

\begin{itemize}
    \item \textbf{Id} -- unikalny identyfikator
    \item \textbf{SubmissionId} -- referencja do artykułu (Submission)
    \item \textbf{ReviewerId} -- referencja do recenzenta (User)
    \item \textbf{AssignedAt} -- data przydzielenia
    \item \textbf{Status} -- status: pending, completed
\end{itemize}

\subsection{Notification (Powiadomienie)}

Encja reprezentująca powiadomienie dla użytkownika:

\begin{itemize}
    \item \textbf{Id} -- unikalny identyfikator
    \item \textbf{UserId} -- referencja do odbiorcy (User)
    \item \textbf{Message} -- treść powiadomienia
    \item \textbf{Type} -- typ: submission, review, decision
    \item \textbf{IsRead} -- flaga czy przeczytane
    \item \textbf{CreatedAt} -- data utworzenia
\end{itemize}

\subsection{FileRecord (Rekord pliku)}

Encja reprezentująca plik przesłany przez użytkownika:

\begin{itemize}
    \item \textbf{Id} -- unikalny identyfikator
    \item \textbf{SubmissionId} -- referencja do artykułu (Submission)
    \item \textbf{FileName} -- nazwa pliku
    \item \textbf{FilePath} -- ścieżka do pliku na serwerze
    \item \textbf{UploadedAt} -- data przesłania
\end{itemize}

\section{Opis implementacji}

\subsection{Backend}

Backend został zbudowany na platformie ASP.NET Core w wersji 8.0 z wykorzystaniem minimal APIs. Architektura backendu składa się z następujących warstw:

\subsubsection{Program.cs -- Punkt wejścia}

Główny plik konfiguracyjny aplikacji ASP.NET Core odpowiada za:

\begin{itemize}
    \item Konfigurację połączenia z bazą danych (SQLite)
    \item Rejestrację usług (DbContext, autentykacja, autoryzacja)
    \item Ustawienie JWT tokenów dla autentykacji
    \item Konfigurację CORS dla komunikacji z frontendem
    \item Inicjalizację bazy danych i dane demo
\end{itemize}

\subsubsection{AppDbContext -- Kontekst bazy danych}

Plik \texttt{Data/AppDbContext.cs} zawiera konfigurację Entity Framework Core:

\begin{itemize}
    \item Definicję DbSets dla wszystkich encji
    \item Konfigurację relacji między encjami
    \item Ustawienia migrac1 bazy danych
\end{itemize}

\subsubsection{Modele}

Folder \texttt{Models/} zawiera definicje wszystkich encji systemu w plikach:

\begin{itemize}
    \item \texttt{User.cs} -- encja użytkownika
    \item \texttt{Submission.cs} -- encja artykułu
    \item \texttt{Review.cs} -- encja recenzji
    \item \texttt{ReviewAssignment.cs} -- encja przydzielenia recenzji
    \item \texttt{Notification.cs} -- encja powiadomienia
    \item \texttt{FileRecord.cs} -- encja rekordu pliku
\end{itemize}

\subsubsection{Endpointy API}

System udostępnia następujące endpointy REST API:

\begin{enumerate}
    \item \textbf{Autentykacja}
    \begin{itemize}
        \item \texttt{POST /api/auth/register} -- rejestracja nowego użytkownika
        \item \texttt{POST /api/auth/login} -- logowanie użytkownika
    \end{itemize}
    
    \item \textbf{Zgłoszenia/Artykuły}
    \begin{itemize}
        \item \texttt{POST /api/submissions} -- przesłanie nowego artykułu
        \item \texttt{GET /api/submissions} -- pobranie własnych artykułów
        \item \texttt{PUT /api/submissions/\{id\}} -- aktualizacja artykułu
        \item \texttt{GET /api/submissions/\{id\}/reviews} -- pobranie recenzji artykułu
    \end{itemize}
    
    \item \textbf{Recenzje}
    \begin{itemize}
        \item \texttt{GET /api/reviews/assigned} -- pobranie przydzielonych recenzji
        \item \texttt{POST /api/submissions/\{id\}/reviews} -- przesłanie recenzji
        \item \texttt{POST /api/submissions/\{id\}/decision} -- podjęcie decyzji redakcyjnej
    \end{itemize}
    
    \item \textbf{Administracja (admin only)}
    \begin{itemize}
        \item \texttt{GET /api/admin/submissions} -- przeglądanie wszystkich artykułów
        \item \texttt{GET /api/admin/reviewers} -- przeglądanie recenzentów
        \item \texttt{POST /api/submissions/\{id\}/assign-reviewer} -- przydzielenie recenzenta
    \end{itemize}
    
    \item \textbf{Powiadomienia}
    \begin{itemize}
        \item \texttt{GET /api/notifications} -- pobranie powiadomień
        \item \texttt{PUT /api/notifications/\{id\}} -- oznaczenie powiadomienia jako przeczytane
    \end{itemize}
\end{enumerate}

\subsection{Frontend}

Frontend został zbudowany na platformie React z narzędziem Vite. Architektura frontendu obejmuje:

\subsubsection{Struktura projektu}

\begin{itemize}
    \item \texttt{src/main.jsx} -- punkt wejścia aplikacji
    \item \texttt{src/App.jsx} -- główny komponent aplikacji
    \item \texttt{src/styles.css} -- style CSS
    \item \texttt{index.html} -- plik HTML
    \item \texttt{vite.config.js} -- konfiguracja Vite
    \item \texttt{package.json} -- zależności i skrypty
\end{itemize}

\subsubsection{Główne komponenty}

\begin{itemize}
    \item \textbf{Login/Register} -- formularz logowania i rejestracji
    \item \textbf{Submission Form} -- formularz do przesyłania artykułów
    \item \textbf{My Submissions} -- lista artykułów przesłanych przez użytkownika
    \item \textbf{Assigned Reviews} -- lista przydzielonych artykułów do recenzowania
    \item \textbf{Notifications} -- panel powiadomień
    \item \textbf{Admin Dashboard} -- panel administracyjny (widoczny tylko dla adminów)
\end{itemize}

\subsubsection{Komunikacja z backendem}

Frontend komunikuje się z backendem za pomocą zapytań HTTP (fetch API):

\begin{itemize}
    \item Autentykacja przez tokeny JWT przechowywane w localStorage
    \item Nagłówek Authorization zawierający token w każdym żądaniu
    \item Obsługa błędów i wyświetlanie komunikatów
    \item Automatyczne odświeżanie danych po zmianach
\end{itemize}

\subsection{Baza danych}

Baza danych SQLite przechowuje wszystkie dane systemu. Struktura bazy jest tworzona automatycznie przez Entity Framework Core na podstawie definicji modeli.

\subsubsection{Tabele}

Tabele w bazie odpowiadają encjom zdefiniowanym w modelu:

\begin{itemize}
    \item \texttt{Users} -- użytkownicy
    \item \texttt{Submissions} -- artykuły
    \item \texttt{Reviews} -- recenzje
    \item \texttt{ReviewAssignments} -- przydzielenia recenzji
    \item \texttt{Notifications} -- powiadomienia
    \item \texttt{FileRecords} -- rekordy plików
\end{itemize}

\subsection{Docker}

Aplikacja jest skonteneryizowana dla łatwego wdrażania. Plik \texttt{docker-compose.yml} definiuje usługi:

\begin{itemize}
    \item Backend -- serwer ASP.NET Core
    \item Frontend -- aplikacja React
    \item SQLite -- baza danych (może być zastąpiona innymi)
\end{itemize}

\subsubsection{Dockerfiles}

\begin{itemize}
    \item \texttt{backend/Dockerfile} -- obraz dla backendu ASP.NET Core
    \item \texttt{frontend/Dockerfile} -- obraz dla frontendu React
\end{itemize}

Obrazy są budowane z wielostopniową budową (multi-stage build) dla optymalizacji rozmiaru.

\section{Stos technologiczny}

\subsection{Backend}

\begin{table}[h]
\centering
\caption{Technologie backendu}
\label{tab:backend_tech}
\begin{tabular}{|l|l|p{4cm}|}
\hline
Technologia & Wersja & Zastosowanie \\
\hline
ASP.NET Core & 8.0 & Framework webowy \\
Entity Framework Core & 7.0.0 & ORM \\
SQLite & - & Baza danych \\
JWT Bearer & 7.0.0 & Autentykacja \\
BCrypt.Net & 4.0.2 & Haszowanie haseł \\
Prometheus & 8.2.1 & Monitoring \\
Swagger & 6.6.1 & Dokumentacja API \\
iText7 & 7.2.5 & Generowanie PDF \\
CsvHelper & 33.0.1 & Obsługa CSV \\
\hline
\end{tabular}
\end{table}

\subsection{Frontend}

\begin{table}[h]
\centering
\caption{Technologie frontendu}
\label{tab:frontend_tech}
\begin{tabular}{|l|l|p{4cm}|}
\hline
Technologia & Wersja & Zastosowanie \\
\hline
React & najnowsza & Framework UI \\
Vite & najnowsza & Build tool \\
JavaScript/JSX & ES2020+ & Język programowania \\
CSS & 3 & Style \\
\hline
\end{tabular}
\end{table}

\section{Konta testowe}

System zawiera predefiniowane konta testowe do celów demonstracyjnych:

\begin{table}[h]
\centering
\caption{Konta testowe}
\label{tab:test_accounts}
\begin{tabular}{|l|l|l|l|}
\hline
Email & Hasło & Rola & Opis \\
\hline
admin@example.com & adminpass & admin & Administrator \\
student@example.com & pass123 & author & Autor artykułów \\
reviewer1@example.com & reviewerpass1 & reviewer & Recenzent 1 \\
reviewer2@example.com & reviewerpass2 & reviewer & Recenzent 2 \\
reviewer3@example.com & reviewerpass3 & reviewer & Recenzent 3 \\
\hline
\end{tabular}
\end{table}

\section{Instalacja i uruchomienie}

\subsection{Wymagania}

\begin{itemize}
    \item .NET 8.0 SDK (dla backendu)
    \item Node.js 18+ (dla frontendu)
    \item Docker i Docker Compose (opcjonalnie)
    \item Git (do klonowania repozytorium)
\end{itemize}

\subsection{Uruchomienie lokalne}

\subsubsection{Backend}

\begin{enumerate}
    \item Przejść do folderu backendu: \texttt{cd backend}
    \item Przywrócić zależności: \texttt{dotnet restore}
    \item Uruchomić aplikację: \texttt{dotnet run}
    \item Backend będzie dostępny na: \texttt{http://localhost:5000}
    \item Dokumentacja Swagger: \texttt{http://localhost:5000/swagger}
\end{enumerate}

\subsubsection{Frontend}

\begin{enumerate}
    \item Przejść do folderu frontendu: \texttt{cd frontend}
    \item Zainstalować zależności: \texttt{npm install}
    \item Uruchomić dev serwer: \texttt{npm run dev}
    \item Frontend będzie dostępny na: \texttt{http://localhost:3000}
\end{enumerate}

\subsubsection{Przygotowanie bazy danych}

Jeśli uruchamiasz aplikację po raz pierwszy:

\begin{enumerate}
    \item Usuń plik \texttt{backend/app.db} (jeśli istnieje)
    \item Uruchom backend -- baza będzie automatycznie utworzona ze schematem
\end{enumerate}

\subsection{Uruchomienie z Docker}

Całą aplikację można uruchomić za pomocą Docker Compose:

\begin{verbatim}
docker compose down
docker compose up --build
\end{verbatim}

To uruchomi backend, frontend i wszystkie wymagane usługi.

\section{Bezpieczeństwo}

\subsection{Autentykacja}

\begin{itemize}
    \item JWT tokeny z sekretem konfigurowalnym przez zmienne środowiskowe
    \item Hasła zahashowane za pomocą BCrypt z odpowiednim salt
    \item Tokeny tracą ważność po określonym czasie
\end{itemize}

\subsection{Autoryzacja}

\begin{itemize}
    \item System ról (admin, author, reviewer)
    \item Polityki autoryzacji dla chronionych endpointów
    \item Kontrola dostępu oparta na rolach (RBAC)
\end{itemize}

\subsection{CORS}

\begin{itemize}
    \item CORS skonfigurowany dla określonych źródeł (localhost)
    \item Możliwość rozszerzenia dla produkcji
\end{itemize}

\subsection{Rekomendacje dla produkcji}

\begin{itemize}
    \item Zamienić SQLite na produkcyjną bazę (PostgreSQL, SQL Server)
    \item Wymusić HTTPS
    \item Skonfigurować certyfikaty SSL/TLS
    \item Zarządzać sekretem JWT w zmiennych środowiskowych
    \item Dodać rate limiting
    \item Włączyć logging i monitoring
    \item Przeprowadzić security audit
\end{itemize}

\section{Przyszłe rozszerzenia}

\subsection{Planowane funkcjonalności}

\begin{itemize}
    \item Obsługa przesyłania plików PDF i innych formatów
    \item Zaawansowana wyszukiwarka artykułów
    \item Komentarze do recenzji
    \item Historia zmian statusu artykułu
    \item Eksport raportów w PDF/CSV
    \item Integracja z email do powiadomień
    \item Dwustopniowa autentykacja (2FA)
    \item Rola moderatora
    \item Automatyczne przypomnienia o terminach
    \item API dla integracji z innymi systemami
\end{itemize}

\subsection{Ulepszenia techniczne}

\begin{itemize}
    \item Migracja na Entity Framework 8+
    \item Dodanie cachingu (Redis)
    \item Implementacja CQRS pattern
    \item Testy jednostkowe i integracyjne
    \item CI/CD pipeline
    \item Dokumentacja API OpenAPI/Swagger
    \item Optymalizacja wydajności bazy danych
    \item Load testing
\end{itemize}

\section{Podsumowanie}

System zgłaszania artykułów naukowych stanowi kompleksową platformę do zarządzania procesem recenzowania artykułów. Architektura oparta na nowoczesnych technologiach webowych zapewnia skalowalność i łatwość rozwijania. Aplikacja jest gotowa do wdrożenia zarówno lokalnie jak i w środowisku produkcyjnym z wykorzystaniem Docker.

Dokumentacja kodu, system roli i autoryzacji, oraz rozdzielenie warstwy prezentacji od logiki biznesowej ułatwia utrzymanie i rozwijanie systemu. Projekt stanowi solidną bazę do dalszego rozwoju i dostosowania do konkretnych potrzeb użytkowników.

\end{document}
